\section{Spectral subspaces under isometric coordinates}

\begin{theorem}[Spectral data after an isometry]
    \label{thm:canonical_isometric_spectral_data}
    \lean{Matrix.exists_canonical_isometric_spectral_data}
    \leanok
    Let $B$ be Hermitian on $\C^\Sigma$, with an orthonormal eigenbasis
    $v_i$ and eigenvalues $\lambda_i$. Put
    $G_c=\operatorname{span}\{v_i:\lambda_i\le c\}$.
    Suppose $G_{\eta_0}=G_{\eta_1}$ and $\|B\|\le U$, where
    $\eta_0,\eta_1,U\in\R$. Let $u:\C^\Sigma\to F$ be a unitary linear
    isometry onto a finite-dimensional complex inner-product space.
    The symmetric operator $K=uBu^{-1}$ has an orthonormal
    eigenbasis $b_j$ and eigenvalues $\kappa_j$ satisfying
    \begin{align}
        \kappa_j & \le\eta_0\quad\text{or}\quad\eta_1\le\kappa_j,\notag \\
        \kappa_j & \le U,\notag                                         \\
        uG_c     & =\operatorname{span}\{b_j:\kappa_j\le c\}.
    \end{align}
    The last equality holds for every $c\in\R$. No order between the
    endpoints or nonzero-dimension assumption is needed.
\end{theorem}
\begin{proof}\leanok
    Any orthonormal eigenbasis spans the same closed low spectral subspace:
    expand an eigenvector in another eigenbasis and compare eigenvalues.
    An isometry carries each eigenspace to the eigenspace of the conjugate
    operator. Apply these facts to an orthonormal spectral basis of $K$.
    Equality of the two supplied low spaces excludes intermediate spectral
    values. Finally $\kappa_j\le|\kappa_j|\le\|K\|=\|B\|\le U$.
    These coordinate assertions support the spectral reduction in
    \cite[Theorem 4.2 and Lemma 2.9]{Arad2016RigorousRG}.
\end{proof}

\begin{theorem}[Two spectral values force a nontrivial local dimension]
    \label{thm:two_spectral_values_local_dimension}
    \lean{Matrix.IsHermitian.two_le_local_dimension_of_two_spectral_values}
    \leanok
    Let $H$ be Hermitian on $(\C^q)^{\otimes L}$. If its spectrum contains
    distinct real values $\mu,\nu$, then $q\ge2$.
\end{theorem}
\begin{proof}\leanok
    If $q\le1$, the configuration set has at most one element. Every
    Hermitian matrix on this space has at most one real spectral value,
    contrary to the two given values. This includes the empty configuration
    set and the empty chain.
\end{proof}
